\documentclass[11pt,a4paper]{article}

% ============================================================
% PACOTES
% ============================================================

\usepackage[utf8]{inputenc}
\usepackage[T1]{fontenc}
\usepackage[portuguese]{babel}
\usepackage{lmodern}

\usepackage{
	amsmath,
	amssymb,
	mathtools
}

\usepackage{siunitx}
\sisetup{
	output-decimal-marker={,},
	per-mode=symbol
}

\usepackage{geometry}
\geometry{
	top=1.65cm,
	bottom=1.7cm,
	left=1.7cm,
	right=1.7cm
}

\usepackage{xcolor}
\usepackage{graphicx}
\usepackage{float}
\usepackage{booktabs}
\usepackage{array}
\usepackage{tabularx}
\usepackage{multicol}
\usepackage{enumitem}

\usepackage{tikz}
\usepackage{pgfplots}

\pgfplotsset{compat=1.18}

\usetikzlibrary{
	arrows.meta,
	positioning,
	calc,
	patterns,
	decorations.pathreplacing,
	shapes.geometric
}

\usepackage[most]{tcolorbox}
\usepackage{fancyhdr}
\usepackage{lastpage}
\usepackage{titlesec}

% ============================================================
% CORES
% ============================================================

\definecolor{eufblue}{RGB}{25,88,145}
\definecolor{euflight}{RGB}{232,243,252}
\definecolor{eufdark}{RGB}{25,48,70}
\definecolor{euforange}{RGB}{230,135,45}
\definecolor{eufgreen}{RGB}{38,127,81}
\definecolor{eufred}{RGB}{190,60,55}
\definecolor{softgray}{RGB}{246,247,249}

% ============================================================
% CONFIGURAÇÃO
% ============================================================

\setlength{\parindent}{0pt}
\setlength{\parskip}{4pt}

\setlist[enumerate]{
	itemsep=5pt,
	topsep=5pt
}

\pagestyle{fancy}
\fancyhf{}

\lhead{\textbf{\textcolor{eufblue}{EUFISICA}}}
\rhead{\textbf{Física e Química A -- 11.º ano}}

\cfoot{
	\small
	Prof. José Gonçalves
	\hfill
	Página \thepage\ de \pageref{LastPage}
}

\setlength{\headheight}{14pt}

\titleformat{\section}
{\large\bfseries\color{eufblue}}
{}
{0pt}
{}

% ============================================================
% CAIXAS
% ============================================================

\newtcolorbox{objetivos}{
	colback=euflight,
	colframe=eufblue,
	boxrule=0.9pt,
	arc=2mm,
	title=\textbf{Objetivos de aprendizagem},
	fonttitle=\bfseries
}

\newtcolorbox{recorda}{
	colback=softgray,
	colframe=eufdark,
	boxrule=0.7pt,
	arc=2mm,
	title=\textbf{Recorda},
	fonttitle=\bfseries
}

\newtcolorbox{atencao}{
	colback=euforange!8,
	colframe=euforange,
	boxrule=0.8pt,
	arc=2mm,
	title=\textbf{Atenção},
	fonttitle=\bfseries
}

\newtcolorbox{desafio}{
	colback=eufgreen!6,
	colframe=eufgreen,
	boxrule=0.8pt,
	arc=2mm,
	title=\textbf{Desafio},
	fonttitle=\bfseries
}

\newtcolorbox{solucao}{
	colback=eufgreen!6,
	colframe=eufgreen!85!black,
	boxrule=0.8pt,
	arc=2mm,
	title=\textbf{Proposta de Resolução},
	fonttitle=\bfseries\small,
	top=4pt,
	bottom=4pt,
	left=6pt,
	right=6pt
}

% ============================================================
% DOCUMENTO
% ============================================================

\begin{document}
	
	% ============================================================
	% CABEÇALHO
	% ============================================================
	
	\begin{center}
		
		{\Huge\bfseries\textcolor{eufblue}{EUFISICA}}
		
		\vspace{1mm}
		
		{\Large\bfseries Ficha de Trabalho --- Proposta de Resolução}
		
		\vspace{1mm}
		
		{\LARGE\bfseries Cinemática}
		
		\vspace{1mm}
		
		{\large
			Gráficos posição--tempo, velocidade--tempo e aceleração
		}
		
		\vspace{2mm}
		
		{\small Prof. José Gonçalves}
		
	\end{center}
	
	\vspace{2mm}
	
	\begin{tabularx}{\textwidth}{|X|X|}
		\hline
		\textbf{Nome:} Solucionário Completo & \textbf{N.º:} --- \\[0.15cm]
		\hline
		\textbf{Turma:} 11.º Ano & \textbf{Data:} 2026/2027 \\[0.15cm]
		\hline
	\end{tabularx}
	
	\vspace{4mm}
	
	% ============================================================
	
	\begin{objetivos}
		
		No final desta ficha deverás ser capaz de:
		
		\begin{itemize}
			
			\item interpretar gráficos posição--tempo, velocidade--tempo e
			aceleração--tempo;
			
			\item determinar a velocidade a partir do declive de um gráfico
			posição--tempo;
			
			\item determinar a aceleração a partir do declive de um gráfico
			velocidade--tempo;
			
			\item determinar deslocamentos através da área sob um gráfico
			velocidade--tempo;
			
			\item distinguir velocidade de rapidez;
			
			\item reconhecer situações de movimento uniforme e movimento uniformemente
			variado;
			
			\item interpretar o sinal das componentes escalares da velocidade e da
			aceleração;
			
			\item representar vetorialmente a velocidade e a aceleração;
			
			\item identificar situações em que um corpo aumenta ou diminui a sua
			rapidez.
			
		\end{itemize}
		
	\end{objetivos}
	
	% ============================================================
	\section*{1. Conceitos fundamentais}
	% ============================================================
	
	\begin{recorda}
		
		Num movimento retilíneo, escolhendo um eixo $Ox$:
		
		\[
		v_x=\frac{\Delta x}{\Delta t}
		\]
		
		e a aceleração média é dada por
		
		\[
		a_x=\frac{\Delta v_x}{\Delta t}.
		\]
		
		Num movimento uniformemente variado:
		
		\[
		v_x=v_{0x}+a_x t
		\]
		
		e
		
		\[
		x=x_0+v_{0x}t+\frac{1}{2}a_xt^2.
		\]
		
		Num gráfico posição--tempo, o \textbf{declive da tangente} corresponde à
		velocidade instantânea.
		
		Num gráfico velocidade--tempo:
		
		\begin{itemize}
			\item o declive corresponde à aceleração;
			\item a área algébrica entre o gráfico e o eixo do tempo corresponde
			ao deslocamento.
		\end{itemize}
		
	\end{recorda}
	
	% ============================================================
	\section*{2. Interpretação de um gráfico posição--tempo}
	% ============================================================
	
	O gráfico representa o movimento retilíneo de um automóvel.
	
	\begin{figure}[H]
		\centering
		
		\begin{tikzpicture}
			
			\begin{axis}[
				width=13cm,
				height=7cm,
				axis lines=left,
				xlabel={Tempo, $t$/s},
				ylabel={Posição, $x$/m},
				xmin=0,
				xmax=14,
				ymin=-25,
				ymax=65,
				xtick={0,2,4,6,8,10,12,14},
				ytick={-20,0,20,40,60},
				grid=both,
				major grid style={gray!30},
				]
				
				\addplot[
				very thick,
				eufblue
				]
				coordinates{
					(0,0)
					(4,40)
					(7,40)
					(11,-20)
					(13,0)
				};
				
				\addplot[
				only marks,
				mark=*,
				mark size=2.5pt,
				eufblue
				]
				coordinates{
					(0,0)
					(4,40)
					(7,40)
					(11,-20)
					(13,0)
				};
				
			\end{axis}
			
		\end{tikzpicture}
		
		\caption{Posição do automóvel em função do tempo.}
		
	\end{figure}
	
	\begin{enumerate}[label=\textbf{\arabic*.}]
		
		\item Indica a posição do automóvel nos instantes:
		
		\begin{enumerate}[label=\alph*)]
			\item $t=0$;
			\item $t=\SI{4}{s}$;
			\item $t=\SI{7}{s}$;
			\item $t=\SI{11}{s}$.
		\end{enumerate}
		
		\begin{solucao}
			\begin{enumerate}[label=\alph*)]
				\item Em $t = \SI{0}{s}$: $x = \SI{0}{m}$.
				\item Em $t = \SI{4}{s}$: $x = \SI{40}{m}$.
				\item Em $t = \SI{7}{s}$: $x = \SI{40}{m}$.
				\item Em $t = \SI{11}{s}$: $x = \SI{-20}{m}$.
			\end{enumerate}
		\end{solucao}
		
		\item Entre que instantes o automóvel esteve em repouso?
		
		\begin{solucao}
			Entre os instantes $t = \SI{4}{s}$ e $t = \SI{7}{s}$ ($t \in [4, 7]\,\text{s}$), pois a posição permanece constante em $x = \SI{40}{m}$ (declive nulo).
		\end{solucao}
		
		\item Em que intervalos de tempo o automóvel se deslocou no sentido positivo
		do eixo $Ox$?
		
		\begin{solucao}
			Nos intervalos em que a posição aumenta (declive positivo): de $t = \SI{0}{s}$ a $t = \SI{4}{s}$ ($[0, 4]\,\text{s}$) e de $t = \SI{11}{s}$ a $t = \SI{13}{s}$ ($[11, 13]\,\text{s}$).
		\end{solucao}
		
		\item Em que intervalo de tempo se deslocou no sentido negativo do eixo
		$Ox$?
		
		\begin{solucao}
			No intervalo de $t = \SI{7}{s}$ a $t = \SI{11}{s}$ ($[7, 11]\,\text{s}$), onde a posição diminui de $\SI{40}{m}$ até $\SI{-20}{m}$ (declive negativo).
		\end{solucao}
		
		\item Determina a componente escalar da velocidade entre
		$t=0$ e $t=\SI{4}{s}$.
		
		\begin{solucao}
			\[
			v_x = \frac{x(\SI{4}{s}) - x(\SI{0}{s})}{4 - 0} = \frac{40 - 0}{4} = \SI{10}{m.s^{-1}}
			\]
		\end{solucao}
		
		\item Determina a componente escalar da velocidade entre
		$t=\SI{7}{s}$ e $t=\SI{11}{s}$.
		
		\begin{solucao}
			\[
			v_x = \frac{x(\SI{11}{s}) - x(\SI{7}{s})}{11 - 7} = \frac{-20 - 40}{4} = \frac{-60}{4} = \SI{-15}{m.s^{-1}}
			\]
		\end{solucao}
		
		\item Explica o significado físico do sinal negativo do resultado anterior.
		
		\begin{solucao}
			O sinal negativo indica que o automóvel se desloca no sentido oposto ao fixado para o eixo $Ox$ (ou seja, no sentido negativo do referencial).
		\end{solucao}
		
		\item Determina a velocidade entre $t=\SI{11}{s}$ e $t=\SI{13}{s}$.
		
		\begin{solucao}
			\[
			v_x = \frac{x(\SI{13}{s}) - x(\SI{11}{s})}{13 - 11} = \frac{0 - (-20)}{2} = \frac{20}{2} = \SI{10}{m.s^{-1}}
			\]
		\end{solucao}
		
		\item Em qual dos intervalos o automóvel apresentou maior rapidez?
		
		Justifica quantitativamente.
		
		\begin{solucao}
			A rapidez corresponde ao módulo da velocidade ($|v_x|$):
			\begin{itemize}
				\item Em $[0, 4]\,\text{s}$: $|v_x| = \SI{10}{m.s^{-1}}$
				\item Em $[4, 7]\,\text{s}$: $|v_x| = \SI{0}{m.s^{-1}}$
				\item Em $[7, 11]\,\text{s}$: $|v_x| = |\SI{-15}{m.s^{-1}}| = \SI{15}{m.s^{-1}}$
				\item Em $[11, 13]\,\text{s}$: $|v_x| = \SI{10}{m.s^{-1}}$
			\end{itemize}
			O automóvel apresentou maior rapidez no intervalo $[7, 11]\,\text{s}$, com um valor de $\SI{15}{m.s^{-1}}$.
		\end{solucao}
		
		\item Determina:
		
		\begin{enumerate}[label=\alph*)]
			\item o deslocamento total;
			\item a distância total percorrida.
		\end{enumerate}
		
		\begin{solucao}
			\begin{enumerate}[label=\alph*)]
				\item $\Delta x = x_{\text{final}} - x_{\text{inicial}} = x(13) - x(0) = 0 - 0 = \SI{0}{m}$.
				\item $s = |40 - 0| + |40 - 40| + |-20 - 40| + |0 - (-20)| = 40 + 0 + 60 + 20 = \SI{120}{m}$.
			\end{enumerate}
		\end{solucao}
		
		\item Explica por que razão o deslocamento e a distância percorrida não são
		iguais neste movimento.
		
		\begin{solucao}
			O deslocamento depende exclusivamente das posições inicial e final ($\Delta x = x_f - x_i$). Como a posição final coincide com a inicial ($x = \SI{0}{m}$), o deslocamento é nulo. Já a distância percorrida ($s$) mede o comprimento efetivo de toda a trajetória percorrida. Houve inversões no sentido do movimento (nos instantes $t = \SI{4}{s}$, $t = \SI{7}{s}$ e $t = \SI{11}{s}$), pelo que $s > |\Delta x|$.
		\end{solucao}
		
	\end{enumerate}
	
	% ============================================================
	\section*{3. Declive e velocidade}
	% ============================================================
	
	Três móveis, A, B e C, deslocam-se ao longo da mesma direção.
	
	\begin{figure}[H]
		\centering
		
		\begin{tikzpicture}
			
			\begin{axis}[
				width=12cm,
				height=7cm,
				axis lines=left,
				xlabel={$t$/s},
				ylabel={$x$/m},
				xmin=0,
				xmax=8,
				ymin=-30,
				ymax=70,
				xtick={0,1,2,3,4,5,6,7,8},
				ytick={-20,0,20,40,60},
				grid=both,
				major grid style={gray!30},
				legend style={
					at={(0.03,0.97)},
					anchor=north west
				}
				]
				
				\addplot[very thick,eufblue]
				coordinates{(0,0)(6,60)};
				\addlegendentry{A}
				
				\addplot[very thick,euforange]
				coordinates{(0,20)(6,50)};
				\addlegendentry{B}
				
				\addplot[very thick,eufred]
				coordinates{(0,40)(6,-20)};
				\addlegendentry{C}
				
			\end{axis}
			
		\end{tikzpicture}
		
		\caption{Gráficos posição--tempo dos móveis A, B e C.}
		
	\end{figure}
	
	\begin{enumerate}[label=\textbf{\arabic*.},resume]
		
		\item Calcula a componente escalar da velocidade de cada móvel.
		
		\begin{solucao}
			\begin{itemize}
				\item Móvel A: $v_A = \frac{60 - 0}{6 - 0} = \SI{10}{m.s^{-1}}$
				\item Móvel B: $v_B = \frac{50 - 20}{6 - 0} = \frac{30}{6} = \SI{5}{m.s^{-1}}$
				\item Móvel C: $v_C = \frac{-20 - 40}{6 - 0} = \frac{-60}{6} = \SI{-10}{m.s^{-1}}$
			\end{itemize}
		\end{solucao}
		
		\item Qual dos móveis apresenta maior rapidez?
		
		\begin{solucao}
			Os móveis \textbf{A} e \textbf{C} apresentam ambos a mesma e maior rapidez, pois $|v_A| = |v_C| = \SI{10}{m.s^{-1}} > |v_B| = \SI{5}{m.s^{-1}}$.
		\end{solucao}
		
		\item Qual dos móveis se desloca no sentido negativo do eixo $Ox$?
		
		\begin{solucao}
			O móvel \textbf{C}, pois a sua velocidade é negativa ($v_C = \SI{-10}{m.s^{-1}} < 0$).
		\end{solucao}
		
		\item Explica a relação entre:
		
		\begin{enumerate}[label=\alph*)]
			\item o sinal do declive e o sentido do movimento;
			\item o valor absoluto do declive e a rapidez.
		\end{enumerate}
		
		\begin{solucao}
			\begin{enumerate}[label=\alph*)]
				\item O sinal do declive indica o sentido do movimento: declive positivo ($>0$) indica movimento no sentido positivo do eixo $Ox$; declive negativo ($<0$) indica movimento no sentido negativo de $Ox$.
				\item O valor absoluto (módulo) do declive do gráfico $x(t)$ é exatamente igual à rapidez do móvel ($|\text{declive}| = |v_x|$). Quanto maior for a inclinação da reta, maior é a rapidez.
			\end{enumerate}
		\end{solucao}
		
	\end{enumerate}
	
	% ============================================================
	\section*{4. Da posição à velocidade}
	% ============================================================
	
	A posição de um carrinho é descrita pela expressão
	
	\[
	x(t)=4+3t,
	\]
	
	em unidades SI.
	
	\begin{enumerate}[label=\textbf{\arabic*.},resume]
		
		\item Qual é a posição inicial do carrinho?
		
		\begin{solucao}
			Comparando com $x(t) = x_0 + v_x t$, a posição inicial é $x_0 = \SI{4}{m}$.
		\end{solucao}
		
		\item Qual é a componente escalar da velocidade?
		
		\begin{solucao}
			A componente escalar da velocidade é $v_x = \SI{3}{m.s^{-1}}$.
		\end{solucao}
		
		\item Em que posição se encontra o carrinho no instante
		$t=\SI{8}{s}$?
		
		\begin{solucao}
			$x(8) = 4 + 3(8) = 4 + 24 = \SI{28}{m}$.
		\end{solucao}
		
		\item Representa, no referencial seguinte, o respetivo gráfico posição--tempo
		para $0\leq t\leq\SI{8}{s}$.
		
		\begin{solucao}
			\begin{center}
				\begin{tikzpicture}
					\begin{axis}[
						width=12cm,
						height=7cm,
						axis lines=left,
						xlabel={$t$/s},
						ylabel={$x$/m},
						xmin=0,
						xmax=9,
						ymin=0,
						ymax=32,
						xtick={0,1,2,3,4,5,6,7,8},
						ytick={0,4,8,12,16,20,24,28,32},
						grid=both,
						major grid style={gray!30}
						]
						\addplot[very thick, eufblue, domain=0:8] {4 + 3*x};
						\addplot[only marks, mark=*, mark size=2.5pt, eufblue] coordinates{(0,4)(8,28)};
					\end{axis}
				\end{tikzpicture}
			\end{center}
		\end{solucao}
		
		\item Representa qualitativamente o gráfico velocidade--tempo correspondente.
		
		\begin{solucao}
			\begin{center}
				\begin{tikzpicture}
					\begin{axis}[
						width=12cm,
						height=5cm,
						axis lines=left,
						xlabel={$t$/s},
						ylabel={$v_x$/m\,s$^{-1}$},
						xmin=0,
						xmax=9,
						ymin=-1,
						ymax=5,
						xtick={0,1,2,3,4,5,6,7,8},
						ytick={-1,0,1,2,3,4,5},
						grid=both,
						major grid style={gray!30}
						]
						\addplot[very thick, eufgreen, domain=0:8] {3};
					\end{axis}
				\end{tikzpicture}
			\end{center}
		\end{solucao}
		
	\end{enumerate}
	
	% ============================================================
	\section*{5. Gráfico velocidade--tempo}
	% ============================================================
	
	O gráfico representa a componente escalar da velocidade de um automóvel.
	
	\begin{figure}[H]
		\centering
		
		\begin{tikzpicture}
			
			\begin{axis}[
				width=13cm,
				height=7cm,
				axis lines=middle,
				xlabel={$t$/s},
				ylabel={$v_x$/m\,s$^{-1}$},
				xmin=0,
				xmax=14,
				ymin=-12,
				ymax=22,
				xtick={0,2,4,6,8,10,12,14},
				ytick={-10,-5,0,5,10,15,20},
				grid=both,
				major grid style={gray!30}
				]
				
				\addplot[
				very thick,
				eufblue
				]
				coordinates{
					(0,0)
					(4,20)
					(7,20)
					(11,0)
					(13,-10)
				};
				
				\addplot[
				only marks,
				mark=*,
				mark size=2.5pt,
				eufblue
				]
				coordinates{
					(0,0)
					(4,20)
					(7,20)
					(11,0)
					(13,-10)
				};
				
			\end{axis}
			
		\end{tikzpicture}
		
		\caption{Componente escalar da velocidade em função do tempo.}
		
	\end{figure}
	
	\begin{enumerate}[label=\textbf{\arabic*.},resume]
		
		\item Qual é a velocidade do automóvel no instante $t=\SI{4}{s}$?
		
		\begin{solucao}
			Em $t = \SI{4}{s}$, a componente escalar da velocidade é $v_x = \SI{20}{m.s^{-1}}$.
		\end{solucao}
		
		\item Entre que instantes o movimento é uniforme?
		
		\begin{solucao}
			Entre os instantes $t = \SI{4}{s}$ e $t = \SI{7}{s}$ ($[4, 7]\,\text{s}$), pois a velocidade se mantém constante em $\SI{20}{m.s^{-1}}$ (aceleração nula).
		\end{solucao}
		
		\item Entre que instantes a rapidez está a aumentar?
		
		\begin{solucao}
			Entre $t = \SI{0}{s}$ e $t = \SI{4}{s}$ ($[0, 4]\,\text{s}$, passa de $\SI{0}{m.s^{-1}}$ a $\SI{20}{m.s^{-1}}$) e entre $t = \SI{11}{s}$ e $t = \SI{13}{s}$ ($[11, 13]\,\text{s}$, o módulo da velocidade passa de $0$ a $\SI{10}{m.s^{-1}}$).
		\end{solucao}
		
		\item Entre que instantes a rapidez está a diminuir?
		
		\begin{solucao}
			Entre $t = \SI{7}{s}$ e $t = \SI{11}{s}$ ($[7, 11]\,\text{s}$), onde a rapidez passa de $\SI{20}{m.s^{-1}}$ até parar em $\SI{0}{m.s^{-1}}$.
		\end{solucao}
		
		\item Em que instante o automóvel inverte o sentido do movimento?
		
		\begin{solucao}
			No instante $t = \SI{11}{s}$, ponto onde a velocidade se anula e muda de sinal (passa de positiva para negativa).
		\end{solucao}
		
		\item Determina a aceleração entre $0$ e \SI{4}{s}.
		
		\begin{solucao}
			\[
			a_x = \frac{v_x(\SI{4}{s}) - v_x(\SI{0}{s})}{4 - 0} = \frac{20 - 0}{4} = \SI{5}{m.s^{-2}}
			\]
		\end{solucao}
		
		\item Determina a aceleração entre \SI{4}{s} e \SI{7}{s}.
		
		\begin{solucao}
			\[
			a_x = \frac{20 - 20}{7 - 4} = \SI{0}{m.s^{-2}}
			\]
		\end{solucao}
		
		\item Determina a aceleração entre \SI{7}{s} e \SI{11}{s}.
		
		\begin{solucao}
			\[
			a_x = \frac{0 - 20}{11 - 7} = \frac{-20}{4} = \SI{-5}{m.s^{-2}}
			\]
		\end{solucao}
		
		\item Determina a aceleração entre \SI{11}{s} e \SI{13}{s}.
		
		\begin{solucao}
			\[
			a_x = \frac{-10 - 0}{13 - 11} = \frac{-10}{2} = \SI{-5}{m.s^{-2}}
			\]
		\end{solucao}
		
		\item Determina o deslocamento do automóvel entre:
		
		\begin{enumerate}[label=\alph*)]
			
			\item $0$ e \SI{4}{s};
			
			\item \SI{4}{s} e \SI{7}{s};
			
			\item \SI{7}{s} e \SI{11}{s};
			
			\item \SI{11}{s} e \SI{13}{s}.
			
		\end{enumerate}
		
		\begin{solucao}
			Calculando geometricamente através das áreas sob o gráfico $v_x(t)$:
			\begin{enumerate}[label=\alph*)]
				\item $\Delta x_1 = \text{Área triângulo} = \frac{4 \times 20}{2} = \SI{40}{m}$.
				\item $\Delta x_2 = \text{Área retângulo} = (7 - 4) \times 20 = 3 \times 20 = \SI{60}{m}$.
				\item $\Delta x_3 = \text{Área triângulo} = \frac{(11 - 7) \times 20}{2} = \frac{4 \times 20}{2} = \SI{40}{m}$.
				\item $\Delta x_4 = \text{Área triângulo (abaixo eixo } t) = \frac{(13 - 11) \times (-10)}{2} = \SI{-10}{m}$.
			\end{enumerate}
		\end{solucao}
		
		\item Determina o deslocamento total entre $t=0$ e $t=\SI{13}{s}$.
		
		\begin{solucao}
			$\Delta x_{\text{total}} = \Delta x_1 + \Delta x_2 + \Delta x_3 + \Delta x_4 = 40 + 60 + 40 + (-10) = \SI{130}{m}$.
		\end{solucao}
		
		\item Determina a distância total percorrida.
		
		\begin{solucao}
			$s = |\Delta x_1| + |\Delta x_2| + |\Delta x_3| + |\Delta x_4| = 40 + 60 + 40 + |-10| = 40 + 60 + 40 + 10 = \SI{150}{m}$.
		\end{solucao}
		
	\end{enumerate}
	
	% ============================================================
	\section*{6. O significado da área num gráfico $v_x(t)$}
	% ============================================================
	
	\begin{recorda}
		
		No gráfico velocidade--tempo,
		
		\[
		\Delta x=\text{área algébrica entre a linha do gráfico e o eixo do tempo}.
		\]
		
		Áreas localizadas acima do eixo do tempo correspondem a deslocamentos
		positivos.
		
		Áreas abaixo do eixo do tempo correspondem a deslocamentos negativos.
		
	\end{recorda}
	
	Considere o gráfico seguinte.
	
	\begin{center}
		
		\begin{tikzpicture}
			
			\begin{axis}[
				width=12cm,
				height=6.5cm,
				axis lines=middle,
				xlabel={$t$/s},
				ylabel={$v_x$/m\,s$^{-1}$},
				xmin=0,
				xmax=10,
				ymin=-8,
				ymax=12,
				xtick={0,2,4,6,8,10},
				ytick={-6,-4,-2,0,2,4,6,8,10},
				grid=both,
				major grid style={gray!30}
				]
				
				\addplot[
				very thick,
				euforange
				]
				coordinates{
					(0,10)
					(4,10)
					(6,0)
					(8,-6)
					(10,-6)
				};
				
			\end{axis}
			
		\end{tikzpicture}
		
	\end{center}
	
	\begin{enumerate}[label=\textbf{\arabic*.},resume]
		
		\item Calcula o deslocamento entre $0$ e \SI{4}{s}.
		
		\begin{solucao}
			$\Delta x_1 = 4 \times 10 = \SI{40}{m}$.
		\end{solucao}
		
		\item Calcula o deslocamento entre \SI{4}{s} e \SI{6}{s}.
		
		\begin{solucao}
			$\Delta x_2 = \frac{(6 - 4) \times 10}{2} = \SI{10}{m}$.
		\end{solucao}
		
		\item Calcula o deslocamento entre \SI{6}{s} e \SI{8}{s}.
		
		\begin{solucao}
			$\Delta x_3 = \frac{(8 - 6) \times (-6)}{2} = \SI{-6}{m}$.
		\end{solucao}
		
		\item Calcula o deslocamento entre \SI{8}{s} e \SI{10}{s}.
		
		\begin{solucao}
			$\Delta x_4 = (10 - 8) \times (-6) = \SI{-12}{m}$.
		\end{solucao}
		
		\item Determina o deslocamento total.
		
		\begin{solucao}
			$\Delta x_{\text{total}} = 40 + 10 + (-6) + (-12) = \SI{32}{m}$.
		\end{solucao}
		
		\item Determina a distância total percorrida.
		
		\begin{solucao}
			$s = 40 + 10 + |-6| + |-12| = 40 + 10 + 6 + 12 = \SI{68}{m}$.
		\end{solucao}
		
		\item Explica a diferença entre os resultados das duas questões anteriores.
		
		\begin{solucao}
			O deslocamento total ($\Delta x = \SI{32}{m}$) é uma grandeza vetorial cuja componente tem em conta o sentido do movimento (a parte percorrida no sentido negativo entra a subtrair). A distância percorrida ($s = \SI{68}{m}$) é uma grandeza escalar sempre positiva que mede a soma de todos os trajetos efetuados.
		\end{solucao}
		
	\end{enumerate}
	
	% ============================================================
	\section*{7. Aceleração}
	% ============================================================
	
	Um automóvel aumenta a sua velocidade de
	\SI{8}{m.s^{-1}} para \SI{20}{m.s^{-1}} durante \SI{6}{s}.
	
	\begin{enumerate}[label=\textbf{\arabic*.},resume]
		
		\item Calcula a aceleração média.
		
		\begin{solucao}
			\[
			a_m = \frac{\Delta v_x}{\Delta t} = \frac{20 - 8}{6} = \frac{12}{6} = \SI{2}{m.s^{-2}}
			\]
		\end{solucao}
		
		\item Interpreta fisicamente o resultado obtido.
		
		\begin{solucao}
			A aceleração de $\SI{2}{m.s^{-2}}$ significa que, em média, a velocidade do automóvel aumenta $\SI{2}{m.s^{-1}}$ a cada segundo que passa.
		\end{solucao}
		
		\item Admitindo aceleração constante, determina a velocidade do automóvel
		\SI{3}{s} depois do início do intervalo considerado.
		
		\begin{solucao}
			Como $a_x = \text{constante}$:
			\[
			v_x(t) = v_{0x} + a_x t = 8 + 2(3) = 8 + 6 = \SI{14}{m.s^{-1}}
			\]
		\end{solucao}
		
	\end{enumerate}
	
	% ============================================================
	\subsection*{7.1. Uma travagem}
	% ============================================================
	
	Um automóvel desloca-se a \SI{25}{m.s^{-1}} quando o condutor inicia uma
	travagem.
	
	Ao fim de \SI{5}{s}, o automóvel encontra-se em repouso.
	
	\begin{enumerate}[label=\textbf{\arabic*.},resume]
		
		\item Determina a aceleração média durante a travagem.
		
		\begin{solucao}
			\[
			a_m = \frac{v_f - v_i}{\Delta t} = \frac{0 - 25}{5} = \SI{-5}{m.s^{-2}}
			\]
		\end{solucao}
		
		\item Por que razão a aceleração apresenta sinal negativo?
		
		\begin{solucao}
			A aceleração tem sinal negativo porque possui sentido oposto ao da velocidade ($v_x > 0$), o que carateriza um movimento retardado (diminuição de rapidez).
		\end{solucao}
		
		\item Admitindo aceleração constante, representa qualitativamente:
		
		\begin{enumerate}[label=\alph*)]
			
			\item o gráfico velocidade--tempo;
			
			\item o gráfico aceleração--tempo.
			
		\end{enumerate}
		
		\begin{solucao}
			\begin{center}
				\begin{minipage}{0.48\textwidth}
					\begin{tikzpicture}
						\begin{axis}[
							width=\linewidth,
							height=5.5cm,
							axis lines=left,
							xlabel={$t$/s},
							ylabel={$v_x$/m\,s$^{-1}$},
							xmin=0,
							xmax=6,
							ymin=0,
							ymax=30,
							grid=both,
							major grid style={gray!30}
							]
							\addplot[very thick, eufblue, domain=0:5] {25 - 5*x};
						\end{axis}
					\end{tikzpicture}
				\end{minipage}
				\hfill
				\begin{minipage}{0.48\textwidth}
					\begin{tikzpicture}
						\begin{axis}[
							width=\linewidth,
							height=5.5cm,
							axis lines=middle,
							xlabel={$t$/s},
							ylabel={$a_x$/m\,s$^{-2}$},
							xmin=0,
							xmax=6,
							ymin=-8,
							ymax=4,
							grid=both,
							major grid style={gray!30}
							]
							\addplot[very thick, eufred, domain=0:5] {-5};
						\end{axis}
					\end{tikzpicture}
				\end{minipage}
			\end{center}
		\end{solucao}
		
		\item Determina a distância percorrida durante a travagem através da área do
		gráfico velocidade--tempo.
		
		\begin{solucao}
			\[
			s = \text{Área do triângulo} = \frac{\text{base} \times \text{altura}}{2} = \frac{5 \times 25}{2} = \SI{62.5}{m}
			\]
		\end{solucao}
		
	\end{enumerate}
	
	% ============================================================
	\section*{8. Velocidade e aceleração como vetores}
	% ============================================================
	
	\begin{recorda}
		
		A velocidade e a aceleração são \textbf{grandezas vetoriais}.
		
		Num movimento retilíneo:
		
		\begin{itemize}
			
			\item se $\vec{v}$ e $\vec{a}$ tiverem o mesmo sentido, a rapidez aumenta;
			
			\item se $\vec{v}$ e $\vec{a}$ tiverem sentidos opostos, a rapidez diminui.
			
		\end{itemize}
		
	\end{recorda}
	
	Analisa as quatro situações seguintes.
	
	% ------------------------------------------------------------
	
	\begin{center}
		
		\begin{tikzpicture}[>=Stealth]
			
			% A
			\node[draw,rounded corners,minimum width=6.3cm,minimum height=2.5cm]
			at (0,0) {};
			
			\node at (-2.5,0.85) {\textbf{Situação A}};
			
			\fill (-1.5,0) circle (0.12);
			
			\draw[->,very thick,eufblue]
			(-1.5,0) -- (0.3,0)
			node[midway,above] {$\vec v$};
			
			\draw[->,very thick,eufred]
			(-1.5,-0.55) -- (-0.3,-0.55)
			node[midway,below] {$\vec a$};
			
			% B
			\node[draw,rounded corners,minimum width=6.3cm,minimum height=2.5cm]
			at (7,0) {};
			
			\node at (4.5,0.85) {\textbf{Situação B}};
			
			\fill (5.5,0) circle (0.12);
			
			\draw[->,very thick,eufblue]
			(5.5,0) -- (7.3,0)
			node[midway,above] {$\vec v$};
			
			\draw[->,very thick,eufred]
			(5.5,-0.55) -- (4.3,-0.55)
			node[midway,below] {$\vec a$};
			
			% C
			\node[draw,rounded corners,minimum width=6.3cm,minimum height=2.5cm]
			at (0,-3.2) {};
			
			\node at (-2.5,-2.35) {\textbf{Situação C}};
			
			\fill (0.7,-3.2) circle (0.12);
			
			\draw[->,very thick,eufblue]
			(0.7,-3.2) -- (-1.1,-3.2)
			node[midway,above] {$\vec v$};
			
			\draw[->,very thick,eufred]
			(0.7,-3.75) -- (-0.5,-3.75)
			node[midway,below] {$\vec a$};
			
			% D
			\node[draw,rounded corners,minimum width=6.3cm,minimum height=2.5cm]
			at (7,-3.2) {};
			
			\node at (4.5,-2.35) {\textbf{Situação D}};
			
			\fill (7.8,-3.2) circle (0.12);
			
			\draw[->,very thick,eufblue]
			(7.8,-3.2) -- (6.0,-3.2)
			node[midway,above] {$\vec v$};
			
			\draw[->,very thick,eufred]
			(7.8,-3.75) -- (9.0,-3.75)
			node[midway,below] {$\vec a$};
			
		\end{tikzpicture}
		
	\end{center}
	
	\begin{enumerate}[label=\textbf{\arabic*.},resume]
		
		\item Em quais situações a rapidez está a aumentar?
		
		\begin{solucao}
			Nas situações \textbf{A} e \textbf{C} ($\vec{v}$ e $\vec{a}$ têm o mesmo sentido).
		\end{solucao}
		
		\item Em quais situações a rapidez está a diminuir?
		
		\begin{solucao}
			Nas situações \textbf{B} e \textbf{D} ($\vec{v}$ e $\vec{a}$ têm sentidos opostos).
		\end{solucao}
		
		\item Na situação A, admitindo o eixo $Ox$ orientado para a direita, indica o
		sinal de:
		
		\begin{enumerate}[label=\alph*)]
			\item $v_x$;
			\item $a_x$.
		\end{enumerate}
		
		\begin{solucao}
			\begin{enumerate}[label=\alph*)]
				\item $v_x > 0$ (positivo)
				\item $a_x > 0$ (positivo)
			\end{enumerate}
		\end{solucao}
		
		\item Repete a questão anterior para as situações B, C e D.
		
		\begin{solucao}
			\begin{itemize}
				\item Situação B: $v_x > 0$ (positivo) e $a_x < 0$ (negativo).
				\item Situação C: $v_x < 0$ (negativo) e $a_x < 0$ (negativo).
				\item Situação D: $v_x < 0$ (negativo) e $a_x > 0$ (positivo).
			\end{itemize}
		\end{solucao}
		
		\item Completa a tabela.
		
		\begin{solucao}
			\begin{center}
				\renewcommand{\arraystretch}{1.3}
				\begin{tabular}{c|c|c|c}
					\toprule
					\textbf{Situação} & \textbf{Sinal de $v_x$} & \textbf{Sinal de $a_x$} & \textbf{Rapidez} \\
					\midrule
					A & $+$ (positivo) & $+$ (positivo) & Aumenta \\
					B & $+$ (positivo) & $-$ (negativo) & Diminui \\
					C & $-$ (negativo) & $-$ (negativo) & Aumenta \\
					D & $-$ (negativo) & $+$ (positivo) & Diminui \\
					\bottomrule
				\end{tabular}
			\end{center}
		\end{solucao}
		
	\end{enumerate}
	
	% ============================================================
	\section*{9. Aceleração positiva nem sempre significa aumento da rapidez}
	% ============================================================
	
	\begin{atencao}
		
		Uma aceleração positiva não significa necessariamente que o corpo esteja a
		aumentar a sua rapidez.
		
		O efeito da aceleração depende simultaneamente dos sinais de $v_x$ e
		$a_x$.
		
	\end{atencao}
	
	Considere quatro automóveis com os seguintes valores instantâneos.
	
	\begin{center}
		
		\renewcommand{\arraystretch}{1.35}
		
		\begin{tabular}{c|c|c}
			
			\toprule
			
			\textbf{Automóvel}
			&
			$v_x$/\si{m.s^{-1}}
			&
			$a_x$/\si{m.s^{-2}}
			
			\\
			
			\midrule
			
			A & $+12$ & $+3$\\
			
			B & $+12$ & $-3$\\
			
			C & $-12$ & $+3$\\
			
			D & $-12$ & $-3$\\
			
			\bottomrule
			
		\end{tabular}
		
	\end{center}
	
	\begin{enumerate}[label=\textbf{\arabic*.},resume]
		
		\item Em quais casos a rapidez está a aumentar?
		
		\begin{solucao}
			Nos casos \textbf{A} e \textbf{D} (onde $v_x$ e $a_x$ têm o mesmo sinal).
		\end{solucao}
		
		\item Em quais casos a rapidez está a diminuir?
		
		\begin{solucao}
			Nos casos \textbf{B} e \textbf{C} (onde $v_x$ e $a_x$ têm sinais opostos).
		\end{solucao}
		
		\item Indica o sentido do movimento de cada automóvel.
		
		\begin{solucao}
			\begin{itemize}
				\item Automóveis \textbf{A} e \textbf{B}: sentido positivo ($v_x > 0$).
				\item Automóveis \textbf{C} e \textbf{D}: sentido negativo ($v_x < 0$).
			\end{itemize}
		\end{solucao}
		
		\item Explica por que razão o automóvel C diminui a rapidez apesar de
		$a_x>0$.
		
		\begin{solucao}
			O automóvel C tem velocidade negativa ($v_x = \SI{-12}{m.s^{-1}}$), ou seja, desloca-se no sentido negativo de $Ox$. Como a aceleração é positiva ($a_x = \SI{+3}{m.s^{-2}}$), o vetor aceleração aponta no sentido positivo. Como $\vec{v}$ e $\vec{a}$ têm sentidos opostos, o movimento é retardado e a rapidez diminui.
		\end{solucao}
		
	\end{enumerate}
	
	% ============================================================
	\section*{10. Movimento uniformemente acelerado}
	% ============================================================
	
	Um carrinho move-se numa trajetória retilínea segundo a equação:
	
	\[
	x(t)=2+4t+t^2
	\]
	
	em unidades SI.
	
	\begin{enumerate}[label=\textbf{\arabic*.},resume]
		
		\item Comparando com
		
		\[
		x=x_0+v_{0x}t+\frac{1}{2}a_xt^2,
		\]
		
		determina:
		
		\begin{enumerate}[label=\alph*)]
			
			\item $x_0$;
			
			\item $v_{0x}$;
			
			\item $a_x$.
			
		\end{enumerate}
		
		\begin{solucao}
			\begin{enumerate}[label=\alph*)]
				\item $x_0 = \SI{2}{m}$
				\item $v_{0x} = \SI{4}{m.s^{-1}}$
				\item $\frac{1}{2}a_x = 1 \implies a_x = \SI{2}{m.s^{-2}}$
			\end{enumerate}
		\end{solucao}
		
		\item Determina a posição do carrinho no instante $t=\SI{4}{s}$.
		
		\begin{solucao}
			$x(4) = 2 + 4(4) + (4)^2 = 2 + 16 + 16 = \SI{34}{m}$.
		\end{solucao}
		
		\item Determina a velocidade no instante $t=\SI{4}{s}$.
		
		\begin{solucao}
			$v_x(t) = v_{0x} + a_x t = 4 + 2t \implies v_x(4) = 4 + 2(4) = \SI{12}{m.s^{-1}}$.
		\end{solucao}
		
		\item Determina o deslocamento entre $t=0$ e $t=\SI{4}{s}$.
		
		\begin{solucao}
			$\Delta x = x(4) - x(0) = 34 - 2 = \SI{32}{m}$.
		\end{solucao}
		
		\item Escreve a equação $v_x(t)$.
		
		\begin{solucao}
			$v_x(t) = 4 + 2t \quad (\text{SI})$.
		\end{solucao}
		
		\item Representa qualitativamente os três gráficos $x(t)$, $v_x(t)$ e
		$a_x(t)$.
		
		\begin{solucao}
			\begin{center}
				\begin{minipage}{0.32\textwidth}
					\centering
					\textbf{$x(t)$}
					
					\begin{tikzpicture}
						\begin{axis}[
							width=\linewidth,
							height=5cm,
							axis lines=left,
							xlabel={$t$},
							ylabel={$x$},
							xmin=0,
							xmax=5,
							ymin=0,
							ymax=45,
							grid=both
							]
							\addplot[very thick, eufblue, domain=0:4] {2 + 4*x + x^2};
						\end{axis}
					\end{tikzpicture}
				\end{minipage}
				\hfill
				\begin{minipage}{0.32\textwidth}
					\centering
					\textbf{$v_x(t)$}
					
					\begin{tikzpicture}
						\begin{axis}[
							width=\linewidth,
							height=5cm,
							axis lines=left,
							xlabel={$t$},
							ylabel={$v_x$},
							xmin=0,
							xmax=5,
							ymin=0,
							ymax=16,
							grid=both
							]
							\addplot[very thick, eufgreen, domain=0:4] {4 + 2*x};
						\end{axis}
					\end{tikzpicture}
				\end{minipage}
				\hfill
				\begin{minipage}{0.32\textwidth}
					\centering
					\textbf{$a_x(t)$}
					
					\begin{tikzpicture}
						\begin{axis}[
							width=\linewidth,
							height=5cm,
							axis lines=left,
							xlabel={$t$},
							ylabel={$a_x$},
							xmin=0,
							xmax=5,
							ymin=0,
							ymax=4,
							grid=both
							]
							\addplot[very thick, eufred, domain=0:4] {2};
						\end{axis}
					\end{tikzpicture}
				\end{minipage}
			\end{center}
		\end{solucao}
		
	\end{enumerate}
	
	% ============================================================
	\section*{11. Relacionar os três gráficos}
	% ============================================================
	
	Um corpo apresenta o seguinte gráfico velocidade--tempo.
	
	\begin{figure}[H]
		\centering
		
		\begin{tikzpicture}
			
			\begin{axis}[
				width=12cm,
				height=6cm,
				axis lines=middle,
				xlabel={$t$/s},
				ylabel={$v_x$/m\,s$^{-1}$},
				xmin=0,
				xmax=10,
				ymin=-8,
				ymax=14,
				xtick={0,2,4,6,8,10},
				ytick={-6,-4,-2,0,2,4,6,8,10,12},
				grid=both,
				major grid style={gray!30}
				]
				
				\addplot[
				very thick,
				eufblue
				]
				coordinates{
					(0,4)
					(4,12)
					(6,12)
					(10,-4)
				};
				
			\end{axis}
			
		\end{tikzpicture}
		
	\end{figure}
	
	\begin{enumerate}[label=\textbf{\arabic*.},resume]
		
		\item Determina a aceleração entre $0$ e \SI{4}{s}.
		
		\begin{solucao}
			$a_x = \frac{12 - 4}{4 - 0} = \frac{8}{4} = \SI{2}{m.s^{-2}}$.
		\end{solucao}
		
		\item Qual é a aceleração entre \SI{4}{s} e \SI{6}{s}?
		
		\begin{solucao}
			$a_x = \SI{0}{m.s^{-2}}$ (movimento uniforme).
		\end{solucao}
		
		\item Determina a aceleração entre \SI{6}{s} e \SI{10}{s}.
		
		\begin{solucao}
			$a_x = \frac{-4 - 12}{10 - 6} = \frac{-16}{4} = \SI{-4}{m.s^{-2}}$.
		\end{solucao}
		
		\item Em que instante, entre \SI{6}{s} e \SI{10}{s}, o corpo fica
		momentaneamente em repouso?
		
		\begin{solucao}
			Determina-se a equação da reta para $t \geq 6$: $v_x(t) = 12 - 4(t - 6)$. \\
			Fazendo $v_x = 0$: $12 - 4(t - 6) = 0 \implies 4(t - 6) = 12 \implies t - 6 = 3 \implies t = \SI{9}{s}$.
		\end{solucao}
		
		\item Em que intervalos:
		
		\begin{enumerate}[label=\alph*)]
			
			\item aumenta a rapidez;
			
			\item diminui a rapidez;
			
			\item apresenta rapidez constante?
			
		\end{enumerate}
		
		\begin{solucao}
			\begin{enumerate}[label=\alph*)]
				\item Aumenta a rapidez em $[0, 4]\,\text{s}$ (passa de $4$ a $\SI{12}{m.s^{-1}}$) e em $[9, 10]\,\text{s}$ (o módulo passa de $0$ a $\SI{4}{m.s^{-1}}$).
				\item Diminui a rapidez em $[6, 9]\,\text{s}$ (passa de $\SI{12}{m.s^{-1}}$ a $0$).
				\item Rapidez constante em $[4, 6]\,\text{s}$ (constante a $\SI{12}{m.s^{-1}}$).
			\end{enumerate}
		\end{solucao}
		
		\item Representa o gráfico aceleração--tempo correspondente.
		
		\begin{solucao}
			\begin{center}
				\begin{tikzpicture}
					\begin{axis}[
						width=12cm,
						height=6cm,
						axis lines=middle,
						xlabel={$t$/s},
						ylabel={$a_x$/m\,s$^{-2}$},
						xmin=0,
						xmax=10,
						ymin=-6,
						ymax=4,
						xtick={0,2,4,6,8,10},
						ytick={-4,-2,0,2},
						grid=both,
						major grid style={gray!30}
						]
						\addplot[very thick, eufred] coordinates{(0,2)(4,2)};
						\addplot[very thick, eufred] coordinates{(4,0)(6,0)};
						\addplot[very thick, eufred] coordinates{(6,-4)(10,-4)};
						\draw[dashed, gray] (4,2) -- (4,0);
						\draw[dashed, gray] (6,0) -- (6,-4);
					\end{axis}
				\end{tikzpicture}
			\end{center}
		\end{solucao}
		
	\end{enumerate}
	
	% ============================================================
	\section*{12. Da tabela para os gráficos}
	% ============================================================
	
	Num laboratório foi registada a velocidade de um carrinho.
	
	\begin{center}
		
		\renewcommand{\arraystretch}{1.35}
		
		\begin{tabular}{c|ccccccc}
			
			\toprule
			
			$t$/s & 0 & 1 & 2 & 3 & 4 & 5 & 6\\
			
			\midrule
			
			$v_x$/\si{m.s^{-1}}
			&
			2 & 4 & 6 & 8 & 10 & 12 & 14
			
			\\
			
			\bottomrule
			
		\end{tabular}
		
	\end{center}
	
	\begin{enumerate}[label=\textbf{\arabic*.},resume]
		
		\item Representa o gráfico velocidade--tempo.
		
		\begin{solucao}
			\begin{center}
				\begin{tikzpicture}
					\begin{axis}[
						width=12cm,
						height=6cm,
						axis lines=left,
						xlabel={$t$/s},
						ylabel={$v_x$/m\,s$^{-1}$},
						xmin=0,
						xmax=7,
						ymin=0,
						ymax=16,
						xtick={0,1,2,3,4,5,6},
						ytick={0,2,4,6,8,10,12,14,16},
						grid=both,
						major grid style={gray!30}
						]
						\addplot[very thick, eufblue, domain=0:6] {2 + 2*x};
						\addplot[only marks, mark=*, mark size=2pt, eufblue] coordinates{
							(0,2)(1,4)(2,6)(3,8)(4,10)(5,12)(6,14)
						};
					\end{axis}
				\end{tikzpicture}
			\end{center}
		\end{solucao}
		
		\item Determina a aceleração do carrinho.
		
		\begin{solucao}
			$a_x = \frac{4 - 2}{1 - 0} = \SI{2}{m.s^{-2}}$.
		\end{solucao}
		
		\item Escreve a equação da velocidade em função do tempo.
		
		\begin{solucao}
			$v_x(t) = 2 + 2t \quad (\text{SI})$.
		\end{solucao}
		
		\item Determina o deslocamento entre $t=0$ e $t=\SI{6}{s}$ utilizando a área
		do gráfico.
		
		\begin{solucao}
			$\Delta x = \text{Área do trapézio} = \frac{(B + b) \times h}{2} = \frac{(14 + 2) \times 6}{2} = \frac{16 \times 6}{2} = \SI{48}{m}$.
		\end{solucao}
		
		\item Admitindo que $x_0=\SI{5}{m}$, determina a posição no instante
		$t=\SI{6}{s}$.
		
		\begin{solucao}
			$x(\SI{6}{s}) = x_0 + \Delta x = 5 + 48 = \SI{53}{m}$.
		\end{solucao}
		
	\end{enumerate}
	
	% ============================================================
	\section*{13. Situação real -- partida de um automóvel}
	% ============================================================
	
	\begin{tcolorbox}[
		colback=euflight,
		colframe=eufblue,
		title=\textbf{Movimento num semáforo}
		]
		
		Um automóvel encontra-se parado num semáforo. Quando a luz passa a verde,
		inicia o movimento numa estrada retilínea.
		
		Durante os primeiros \SI{6}{s}, a sua aceleração é aproximadamente
		constante e igual a \SI{2.5}{m.s^{-2}}.
		
		Depois mantém a velocidade atingida durante \SI{10}{s}.
		
		Finalmente, o condutor trava uniformemente durante \SI{5}{s}, até o
		automóvel parar.
		
	\end{tcolorbox}
	
	\begin{enumerate}[label=\textbf{\arabic*.},resume]
		
		\item Qual é a velocidade ao fim dos primeiros \SI{6}{s}?
		
		\begin{solucao}
			$v_x(6) = 0 + 2,5 \times 6 = \SI{15}{m.s^{-1}}$.
		\end{solucao}
		
		\item Representa qualitativamente o gráfico velocidade--tempo de todo o
		movimento.
		
		\begin{solucao}
			\begin{center}
				\begin{tikzpicture}
					\begin{axis}[
						width=13cm,
						height=7cm,
						axis lines=left,
						xlabel={$t$/s},
						ylabel={$v_x$/m\,s$^{-1}$},
						xmin=0,
						xmax=23,
						ymin=0,
						ymax=18,
						xtick={0,2,4,6,8,10,12,14,16,18,20,21},
						ytick={0,3,6,9,12,15,18},
						grid=both,
						major grid style={gray!30}
						]
						\addplot[very thick, eufblue] coordinates{(0,0)(6,15)(16,15)(21,0)};
					\end{axis}
				\end{tikzpicture}
			\end{center}
		\end{solucao}
		
		\item Determina a aceleração durante a travagem.
		
		\begin{solucao}
			A travagem ocorre entre os instantes $t = \SI{16}{s}$ e $t = \SI{21}{s}$ ($\Delta t = \SI{5}{s}$):
			\[
			a_x = \frac{0 - 15}{21 - 16} = \frac{-15}{5} = \SI{-3}{m.s^{-2}}
			\]
		\end{solucao}
		
		\item Determina a distância percorrida durante os primeiros \SI{6}{s}.
		
		\begin{solucao}
			$s_1 = \frac{6 \times 15}{2} = \SI{45}{m}$.
		\end{solucao}
		
		\item Determina a distância percorrida durante os \SI{10}{s} de movimento
		uniforme.
		
		\begin{solucao}
			$s_2 = 10 \times 15 = \SI{150}{m}$.
		\end{solucao}
		
		\item Determina a distância percorrida durante a travagem.
		
		\begin{solucao}
			$s_3 = \frac{5 \times 15}{2} = \SI{37.5}{m}$.
		\end{solucao}
		
		\item Determina a distância total percorrida.
		
		\begin{solucao}
			$s_{\text{total}} = s_1 + s_2 + s_3 = 45 + 150 + 37,5 = \SI{232.5}{m}$.
		\end{solucao}
		
		\item Calcula a velocidade média durante todo o percurso.
		
		\begin{solucao}
			\[
			v_m = \frac{\Delta x_{\text{total}}}{\Delta t_{\text{total}}} = \frac{232,5}{21} \approx \SI{11.07}{m.s^{-1}}
			\]
		\end{solucao}
		
	\end{enumerate}
	
	% ============================================================
	\section*{14. Problema de síntese}
	% ============================================================
	
	\begin{desafio}
		
		Um móvel desloca-se numa trajetória retilínea.
		
		No instante $t=0$ encontra-se na posição $x_0=\SI{10}{m}$ e apresenta
		velocidade
		
		\[
		v_{0x}=-\SI{8}{m.s^{-1}}.
		\]
		
		A aceleração é constante e igual a
		
		\[
		a_x=+\SI{2}{m.s^{-2}}.
		\]
		
	\end{desafio}
	
	\begin{enumerate}[label=\textbf{\arabic*.},resume]
		
		\item Representa os vetores velocidade e aceleração no instante inicial,
		admitindo o eixo $Ox$ orientado para a direita.
		
		\begin{solucao}
			\begin{center}
				\begin{tikzpicture}[>=Stealth]
					\draw[->,thick] (-5,0)--(5,0) node[right] {$+x$};
					\fill (0,0) circle (0.12);
					\node[below] at (0,-0.2) {móvel};
					\draw[->,very thick,eufblue] (0,0.4) -- (-2,0.4) node[midway,above] {$\vec{v}_0$};
					\draw[->,very thick,eufred] (0,0.8) -- (1.5,0.8) node[midway,above] {$\vec{a}$};
				\end{tikzpicture}
			\end{center}
		\end{solucao}
		
		\item No instante inicial, a rapidez está a aumentar ou a diminuir?
		
		Justifica.
		
		\begin{solucao}
			A rapidez está a \textbf{diminuir}. Justificação: os vetores velocidade ($\vec{v}_0$) e aceleração ($\vec{a}$) têm sentidos opostos ($v_{0x} < 0$ e $a_x > 0$), o que carateriza um movimento retardado.
		\end{solucao}
		
		\item Escreve a equação da velocidade:
		
		\begin{solucao}
			\[
			v_x(t) = -8 + 2t \quad (\text{SI})
			\]
		\end{solucao}
		
		\item Em que instante o móvel fica momentaneamente em repouso?
		
		\begin{solucao}
			Fazendo $v_x(t) = 0 \implies -8 + 2t = 0 \implies 2t = 8 \implies t = \SI{4}{s}$.
		\end{solucao}
		
		\item Que alteração ocorre no sentido do movimento nesse instante?
		
		\begin{solucao}
			O móvel \textbf{inverte o sentido do seu movimento}: até $t = \SI{4}{s}$ deslocava-se no sentido negativo ($v_x < 0$) e, a partir desse instante, passa a deslocar-se no sentido positivo ($v_x > 0$).
		\end{solucao}
		
		\item Determina a posição nesse instante.
		
		\begin{solucao}
			$x(4) = 10 - 8(4) + (4)^2 = 10 - 32 + 16 = \SI{-6}{m}$.
		\end{solucao}
		
		\item Escreve a equação da posição em função do tempo.
		
		\begin{solucao}
			\[
			x(t) = 10 - 8t + t^2 \quad (\text{SI})
			\]
		\end{solucao}
		
		\item Determina a posição no instante $t=\SI{8}{s}$.
		
		\begin{solucao}
			$x(8) = 10 - 8(8) + 8^2 = 10 - 64 + 64 = \SI{10}{m}$.
		\end{solucao}
		
		\item Determina o deslocamento entre $t=0$ e $t=\SI{8}{s}$.
		
		\begin{solucao}
			$\Delta x = x(8) - x(0) = 10 - 10 = \SI{0}{m}$.
		\end{solucao}
		
		\item Determina a distância total percorrida durante esses \SI{8}{s}.
		
		\begin{solucao}
			\begin{itemize}
				\item De $t = \SI{0}{s}$ a $t = \SI{4}{s}$: o móvel vai de $x = \SI{10}{m}$ a $x = \SI{-6}{m} \implies s_1 = |-6 - 10| = \SI{16}{m}$.
				\item De $t = \SI{4}{s}$ a $t = \SI{8}{s}$: o móvel vai de $x = \SI{-6}{m}$ a $x = \SI{10}{m} \implies s_2 = |10 - (-6)| = \SI{16}{m}$.
			\end{itemize}
			Distância total: $s = s_1 + s_2 = 16 + 16 = \SI{32}{m}$.
		\end{solucao}
		
		\item Esboça o gráfico velocidade--tempo entre $0$ e \SI{8}{s}.
		
		\begin{solucao}
			\begin{center}
				\begin{tikzpicture}
					\begin{axis}[
						width=12cm,
						height=6cm,
						axis lines=middle,
						xlabel={$t$/s},
						ylabel={$v_x$/m\,s$^{-1}$},
						xmin=0,
						xmax=9,
						ymin=-10,
						ymax=10,
						xtick={0,1,2,3,4,5,6,7,8},
						ytick={-10,-8,-6,-4,-2,0,2,4,6,8,10},
						grid=both,
						major grid style={gray!30}
						]
						\addplot[very thick, eufblue, domain=0:8] {-8 + 2*x};
					\end{axis}
				\end{tikzpicture}
			\end{center}
		\end{solucao}
		
		\item Esboça o gráfico aceleração--tempo.
		
		\begin{solucao}
			\begin{center}
				\begin{tikzpicture}
					\begin{axis}[
						width=12cm,
						height=5.5cm,
						axis lines=middle,
						xlabel={$t$/s},
						ylabel={$a_x$/m\,s$^{-2}$},
						xmin=0,
						xmax=9,
						ymin=-3,
						ymax=4,
						xtick={0,1,2,3,4,5,6,7,8},
						ytick={-3,-2,-1,0,1,2,3,4},
						grid=both,
						major grid style={gray!30}
						]
						\addplot[very thick, eufred, domain=0:8] {2};
					\end{axis}
				\end{tikzpicture}
			\end{center}
		\end{solucao}
		
		\item Indica qualitativamente a forma do gráfico posição--tempo.
		
		É uma reta ou uma curva? Justifica.
		
		\begin{solucao}
			É uma \textbf{curva} (especificamente uma \textbf{parábola} com concavidade voltada para cima, uma vez que $a_x = \SI{+2}{m.s^{-2}} > 0$). Justificação: a posição varia de forma quadrática com o tempo ($x \propto t^2$), visto tratar-se de um movimento uniformemente variado com aceleração constante não nula.
		\end{solucao}
		
	\end{enumerate}
	
	% ============================================================
	\section*{15. Questões de escolha múltipla}
	% ============================================================
	
	\begin{enumerate}[label=\textbf{\arabic*.},resume]
		
		\item Num gráfico posição--tempo, uma reta com declive negativo significa que
		o corpo:
		
		\begin{enumerate}[label=\Alph*.]
			\item está necessariamente a diminuir a rapidez;
			\item está em repouso;
			\item \textbf{desloca-se no sentido negativo do referencial;}
			\item apresenta aceleração negativa.
		\end{enumerate}
		
		\begin{solucao}
			\textbf{Opção C.} O declive do gráfico $x(t)$ dá a velocidade. Declive negativo significa $v_x < 0$, ou seja, movimento no sentido negativo.
		\end{solucao}
		
		\item Num gráfico velocidade--tempo, uma linha horizontal acima do eixo do
		tempo representa:
		
		\begin{enumerate}[label=\Alph*.]
			\item \textbf{velocidade positiva constante;}
			\item aceleração positiva constante;
			\item repouso;
			\item velocidade crescente.
		\end{enumerate}
		
		\begin{solucao}
			\textbf{Opção A.} Reta horizontal em $v_x > 0$ significa velocidade positiva que não muda com o tempo (MRU).
		\end{solucao}
		
		\item Se $v_x<0$ e $a_x>0$, então:
		
		\begin{enumerate}[label=\Alph*.]
			\item a rapidez aumenta sempre;
			\item \textbf{a rapidez diminui enquanto o corpo mantiver $v_x<0$;}
			\item o corpo está em repouso;
			\item a aceleração é nula.
		\end{enumerate}
		
		\begin{solucao}
			\textbf{Opção B.} Sinais opostos em $v_x$ e $a_x$ significam que $\vec{v}$ e $\vec{a}$ têm sentidos contrários, reduzindo a rapidez.
		\end{solucao}
		
		\item A área algébrica sob um gráfico velocidade--tempo corresponde:
		
		\begin{enumerate}[label=\Alph*.]
			\item à aceleração;
			\item à rapidez;
			\item \textbf{ao deslocamento;}
			\item à posição inicial.
		\end{enumerate}
		
		\begin{solucao}
			\textbf{Opção C.} A área algébrica sob o gráfico $v_x(t)$ é numericamente igual ao deslocamento ($\Delta x$).
		\end{solucao}
		
		\item Um automóvel apresenta $v_x=+\SI{15}{m.s^{-1}}$ e
		$a_x=-\SI{3}{m.s^{-2}}$.
		
		Nesse instante:
		
		\begin{enumerate}[label=\Alph*.]
			\item desloca-se no sentido positivo e aumenta a rapidez;
			\item desloca-se no sentido negativo e aumenta a rapidez;
			\item \textbf{desloca-se no sentido positivo e diminui a rapidez;}
			\item desloca-se no sentido negativo e diminui a rapidez.
		\end{enumerate}
		
		\begin{solucao}
			\textbf{Opção C.} Como $v_x > 0$, desloca-se no sentido positivo. Como $v_x$ e $a_x$ têm sinais contrários, a rapidez diminui.
		\end{solucao}
		
	\end{enumerate}
	
	% ============================================================
	\section*{16. Verdadeiro ou falso}
	% ============================================================
	
	Classifica cada afirmação como verdadeira (V) ou falsa (F).
	Corrige as falsas.
	
	\begin{enumerate}[label=\textbf{\arabic*.},resume]
		
		\item[\textbf{[V]}]
		Num gráfico posição--tempo, a inclinação permite determinar a velocidade.
		
		\begin{solucao}
			\textbf{Verdadeiro.} O declive da reta (ou tangente) em $x(t)$ é igual à velocidade instantânea.
		\end{solucao}
		
		\item[\textbf{[F]}]
		Uma velocidade negativa significa obrigatoriamente que o corpo está a
		abrandar.
		
		\begin{solucao}
			\textbf{Falso.} Significa apenas que o corpo se move no sentido negativo do eixo $Ox$. Se a aceleração também for negativa ($a_x < 0$), o corpo estará a aumentar a sua rapidez.
		\end{solucao}
		
		\item[\textbf{[V]}]
		Uma aceleração nula implica velocidade constante.
		
		\begin{solucao}
			\textbf{Verdadeiro.} Se $a_x = 0$, não há variação da velocidade no tempo, tratando-se de um Movimento Retilíneo Uniforme.
		\end{solucao}
		
		\item[\textbf{[F]}]
		Uma aceleração negativa implica sempre diminuição da rapidez.
		
		\begin{solucao}
			\textbf{Falso.} Se a velocidade também for negativa ($v_x < 0$), uma aceleração negativa faz aumentar a rapidez (movimento acelerado no sentido negativo).
		\end{solucao}
		
		\item[\textbf{[V]}]
		Se velocidade e aceleração tiverem o mesmo sentido, a rapidez aumenta.
		
		\begin{solucao}
			\textbf{Verdadeiro.} Vetores com a mesma direção e sentido resultam em movimento acelerado.
		\end{solucao}
		
		\item[\textbf{[V]}]
		A área sob um gráfico velocidade--tempo pode assumir um valor negativo.
		
		\begin{solucao}
			\textbf{Verdadeiro.} Áreas situadas abaixo do eixo dos tempos $t$ entram com sinal negativo na área algébrica, representando deslocamentos no sentido negativo.
		\end{solucao}
		
		\item[\textbf{[V]}]
		Num gráfico velocidade--tempo, uma reta inclinada corresponde a aceleração
		não nula.
		
		\begin{solucao}
			\textbf{Verdadeiro.} O declive de $v_x(t)$ representa a aceleração; se a reta for inclinada, o declive é diferente de zero.
		\end{solucao}
		
	\end{enumerate}
	
	% ============================================================
	\section*{17. Desafio final -- identifica o movimento}
	% ============================================================
	
	São apresentados quatro gráficos velocidade--tempo.
	
	\begin{center}
		
		\begin{tikzpicture}
			
			% gráfico A
			\begin{axis}[
				at={(0cm,6cm)},
				anchor=south west,
				width=6cm,
				height=5cm,
				axis lines=middle,
				xmin=0,
				xmax=5,
				ymin=-1,
				ymax=6,
				xlabel={$t$},
				ylabel={$v_x$},
				title={A},
				grid=both
				]
				\addplot[very thick,eufblue]
				coordinates{(0,2)(5,2)};
			\end{axis}
			
			% gráfico B
			\begin{axis}[
				at={(7cm,6cm)},
				anchor=south west,
				width=6cm,
				height=5cm,
				axis lines=middle,
				xmin=0,
				xmax=5,
				ymin=-1,
				ymax=6,
				xlabel={$t$},
				ylabel={$v_x$},
				title={B},
				grid=both
				]
				\addplot[very thick,eufgreen]
				coordinates{(0,0)(5,5)};
			\end{axis}
			
			% gráfico C
			\begin{axis}[
				at={(0cm,0cm)},
				anchor=south west,
				width=6cm,
				height=5cm,
				axis lines=middle,
				xmin=0,
				xmax=5,
				ymin=-1,
				ymax=6,
				xlabel={$t$},
				ylabel={$v_x$},
				title={C},
				grid=both
				]
				\addplot[very thick,euforange]
				coordinates{(0,5)(5,0)};
			\end{axis}
			
			% gráfico D
			\begin{axis}[
				at={(7cm,0cm)},
				anchor=south west,
				width=6cm,
				height=5cm,
				axis lines=middle,
				xmin=0,
				xmax=5,
				ymin=-6,
				ymax=2,
				xlabel={$t$},
				ylabel={$v_x$},
				title={D},
				grid=both
				]
				\addplot[very thick,eufred]
				coordinates{(0,-1)(5,-5)};
			\end{axis}
			
		\end{tikzpicture}
		
	\end{center}
	
	\vspace{3mm}
	
	\begin{solucao}
		\textbf{Associação de descrições e gráficos:}
		\begin{itemize}
			\item \textbf{Gráfico A:} Movimento uniforme no sentido positivo.
			\item \textbf{Gráfico B:} Movimento no sentido positivo com aumento da rapidez.
			\item \textbf{Gráfico C:} Movimento no sentido positivo com diminuição da rapidez.
			\item \textbf{Gráfico D:} Movimento no sentido negativo com aumento da rapidez.
		\end{itemize}
		
		\textbf{Explicação dos sinais de $v_x$ e $a_x$:}
		\begin{itemize}
			\item \textbf{Gráfico A:} $v_x > 0$ (linha acima do eixo $t$) e $a_x = 0$ (declive nulo). Velocidade positiva e constante.
			\item \textbf{Gráfico B:} $v_x > 0$ e $a_x > 0$ (declive positivo). Sinais iguais $\implies$ movimento acelerado no sentido positivo.
			\item \textbf{Gráfico C:} $v_x > 0$ e $a_x < 0$ (declive negativo). Sinais opostos $\implies$ movimento retardado no sentido positivo.
			\item \textbf{Gráfico D:} $v_x < 0$ (linha abaixo do eixo $t$) e $a_x < 0$ (declive negativo). Sinais iguais $\implies$ movimento acelerado no sentido negativo (módulo da velocidade aumenta).
		\end{itemize}
	\end{solucao}
	
	% ============================================================
	\section*{18. Autoavaliação}
	% ============================================================
	
	\begin{center}
		
		\renewcommand{\arraystretch}{1.45}
		
		\begin{tabularx}{\textwidth}{Xccc}
			
			\toprule
			
			\textbf{Consigo...}
			&
			\textbf{Sim}
			&
			\textbf{Com ajuda}
			&
			\textbf{Ainda não}
			
			\\
			
			\midrule
			
			Interpretar gráficos posição--tempo
			&
			$\blacksquare$
			&
			$\square$
			&
			$\square$
			\\
			
			Calcular velocidades a partir de gráficos $x(t)$
			&
			$\blacksquare$
			&
			$\square$
			&
			$\square$
			\\
			
			Interpretar gráficos velocidade--tempo
			&
			$\blacksquare$
			&
			$\square$
			&
			$\square$
			\\
			
			Determinar deslocamentos através de áreas
			&
			$\blacksquare$
			&
			$\square$
			&
			$\square$
			\\
			
			Calcular acelerações
			&
			$\blacksquare$
			&
			$\square$
			&
			$\square$
			\\
			
			Relacionar os sinais de $v_x$ e $a_x$
			&
			$\blacksquare$
			&
			$\square$
			&
			$\square$
			\\
			
			Representar vetores velocidade e aceleração
			&
			$\blacksquare$
			&
			$\square$
			&
			$\square$
			\\
			
			Relacionar os gráficos $x(t)$, $v(t)$ e $a(t)$
			&
			$\blacksquare$
			&
			$\square$
			&
			$\square$
			\\
			
			\bottomrule
			
		\end{tabularx}
		
	\end{center}
	
	\vfill
	
	% ============================================================
	% RODAPÉ FINAL
	% ============================================================
	
	\begin{center}
		
		\begin{tikzpicture}
			
			\node[
			rounded corners=3mm,
			draw=eufblue,
			line width=1pt,
			fill=euflight,
			inner xsep=10mm,
			inner ysep=4mm
			]
			{
				\begin{minipage}{0.84\textwidth}
					
					\centering
					
					{\Large\bfseries\textcolor{eufblue}{EUFISICA}}
					
					\vspace{1mm}
					
					\small
					
					Posição $\bullet$
					Velocidade $\bullet$
					Aceleração $\bullet$
					Gráficos $\bullet$
					Vetores
					
					\vspace{1mm}
					
					\textit{
						Compreender um movimento é conseguir contar a mesma história
						através de equações, vetores e gráficos.
					}
					
					\vspace{1mm}
					
					\textbf{Prof. José Gonçalves}
					
				\end{minipage}
			};
			
		\end{tikzpicture}
		
	\end{center}
	
\end{document}